\documentclass[12pt,a4paper]{article}

\usepackage[utf8]{inputenc}
\usepackage[T1]{fontenc}
\usepackage[spanish]{babel}
\usepackage[left=2.5cm, right=2.5cm, top=2.5cm, bottom=2.5cm]{geometry}
\usepackage{graphicx}
\usepackage{booktabs}
\usepackage{array}
\usepackage{longtable}
\usepackage{amsmath}
\usepackage{xcolor}
\usepackage{hyperref}
\usepackage{float}
\usepackage{caption}
\usepackage{subcaption}
\usepackage{enumitem}
\usepackage{fancyhdr}
\usepackage{titlesec}
\usepackage{parskip}
\usepackage{mdframed}
\usepackage{lmodern}

\hypersetup{
    colorlinks=true,
    linkcolor=black,
    urlcolor=blue,
    citecolor=black
}

\pagestyle{fancy}
\fancyhf{}
\rhead{Sprint 2 -- Modelización y Política de Rotación}
\lhead{Proyecto SAMO -- Sostenibilidad y Ciencia}
\rfoot{Página \thepage}
\lfoot{Equipo Chacorocalfarobar}
\renewcommand{\headrulewidth}{0.4pt}
\renewcommand{\footrulewidth}{0.4pt}

\titleformat{\section}{\large\bfseries\color{black}}{}{0em}{\thesection\quad}[\titlerule]
\titleformat{\subsection}{\normalsize\bfseries}{}{0em}{\thesubsection\quad}
\titleformat{\subsubsection}{\normalsize\itshape}{}{0em}{\thesubsubsection\quad}

\definecolor{graybox}{gray}{0.93}
\definecolor{accentblue}{RGB}{30, 90, 160}

\newmdenv[
  backgroundcolor=graybox,
  linecolor=accentblue,
  linewidth=1.2pt,
  leftline=true,
  rightline=false,
  topline=false,
  bottomline=false,
  innerleftmargin=10pt,
  innerrightmargin=8pt,
  innertopmargin=6pt,
  innerbottommargin=6pt
]{reglabox}

\begin{document}

% ─── PORTADA ────────────────────────────────────────────────────────────────
\begin{titlepage}
    \centering
    \vspace*{1.5cm}
    {\Large \textbf{Sostenibilidad y Ciencia}}\\[0.4cm]
    {\large Proyecto: Análisis y Optimización Operativa de Sistemas Agrovoltaicos Dinámicos}\\[2.5cm]

    {\Huge \textbf{E03 -- Informe de Modelización}}\\[0.4cm]
    {\LARGE \textbf{Política de Rotación de Paneles Agrovoltaicos}}\\[0.4cm]
    {\large Sprint 2}\\[2cm]

    \rule{10cm}{0.5pt}\\[1cm]

    \begin{tabular}{rl}
        \textbf{Proyecto:}   & SAMO -- Sistemas Agrovoltaicos Dinámicos \\[4pt]
        \textbf{Equipo:}     & Chacorocalfarobar \\[4pt]
        \textbf{Director PM:}& Daniel Álvarez Sarroca \\[4pt]
        \textbf{Sprint:}     & 2 de 4 \\[4pt]
        \textbf{Entregable:} & E03 \\[4pt]
        \textbf{Fecha:}      & Abril 2026 \\[4pt]
        \textbf{Versión:}    & 1.0
    \end{tabular}

    \vfill
    {\small Documento técnico interno -- Pendiente de validación agronómica con el especialista de planta}
\end{titlepage}

% ─── RESUMEN EJECUTIVO ──────────────────────────────────────────────────────
\newpage
\section*{Resumen Ejecutivo}
\addcontentsline{toc}{section}{Resumen Ejecutivo}

Este informe documenta el trabajo técnico realizado durante el Sprint~2 del proyecto SAMO.
El objetivo del sprint era desarrollar un modelo de decisión que permitiera proponer una
primera política operativa de rotación de paneles agrovoltaicos, equilibrando
simultáneamente la producción energética y las condiciones microclimáticas del cultivo.

Se integraron ocho conjuntos de datos de sensores a una granularidad temporal común de
seis horas, obteniendo 337 observaciones modelables correspondientes al período
junio--octubre de 2025. Se definió un índice combinado Energía--Cultivo (\textbf{IEC})
como variable objetivo provisional, y se entrenaron dos modelos complementarios:
\textit{Decision Tree Regressor} para la extracción de reglas interpretables y
\textit{ElasticNetCV} para la generalización predictiva.

Los resultados más relevantes son:

\begin{itemize}[noitemsep]
    \item El modo de operación \textbf{TRACKING\_PM} obtiene un IEC medio de \textbf{0.742},
          frente a \textbf{0.204} del modo HORIZONTAL.
    \item Se han derivado \textbf{6 reglas candidatas de rotación}, trazables y revisables
          por el equipo técnico y agrícola.
    \item El modelo ElasticNet alcanza un R$^2$ de \textbf{0.873} en el conjunto de test,
          con un MAE de 0.057.
    \item Se identifican condiciones críticas de alta y baja prioridad basadas en Albedo,
          elevación solar y humedad del suelo.
\end{itemize}

Todas las reglas se consideran \textit{candidatas observacionales} pendientes de validación
con el especialista agrícola antes de integrarse en el dashboard del Sprint~3.

\vspace{0.5cm}
\tableofcontents

% ─── 1. INTRODUCCIÓN ────────────────────────────────────────────────────────
\newpage
\section{Introducción}

\subsection{Contexto del proyecto}

Los sistemas agrovoltaicos dinámicos combinan la generación de energía fotovoltaica con
el cultivo agrícola bajo los propios paneles solares. A diferencia de las instalaciones
estáticas, los paneles disponen de un eje de rotación que permite modificar su ángulo de
inclinación a lo largo del día. Esta capacidad introduce un grado de libertad operativo
fundamental: el ángulo de los paneles afecta simultáneamente a la irradiancia recibida
por el sistema fotovoltaico y al microclima del cultivo situado debajo.

El objetivo del proyecto no es maximizar únicamente la energía producida, sino encontrar
un \textbf{equilibrio dinámico} entre producción energética y condiciones favorables para
el cultivo. Esto incluye variables como la radiación fotosintéticamente activa (ePAR),
el contenido volumétrico de agua del suelo (VWC), la temperatura del suelo y el estrés
térmico del cultivo.

\subsection{Objetivos del Sprint 2}

El Sprint~2 tiene como misión principal transformar los hallazgos exploratorios del
Sprint~1 en un modelo de decisión operativo. De forma específica, se busca:

\begin{enumerate}[noitemsep]
    \item Definir un target cuantificable que capture el equilibrio energía--cultivo.
    \item Preparar un dataset de modelización limpio y reproducible.
    \item Entrenar y evaluar al menos dos enfoques de modelización complementarios.
    \item Traducir los resultados del modelo a reglas operativas interpretables.
    \item Publicar el código en formato notebook replicable.
\end{enumerate}

\subsection{Alcance del informe}

Este documento cubre los entregables técnicos \textbf{E03} (informe de política de
rotación) y está vinculado al entregable \textbf{E04} (código del modelo, notebook
\texttt{sprint2\_modelizacion\_agrovoltaica.ipynb}). No se incluyen en este informe los
artefactos de gestión del sprint (actas, presupuesto, registro de riesgos).

% ─── 2. PIPELINE DE DATOS ───────────────────────────────────────────────────
\newpage
\section{Pipeline de Datos}

\subsection{Datasets integrados}

Se han procesado ocho conjuntos de datos en formato CSV procedentes de los sensores de
la planta agrovoltaica. La integración se realiza mediante la variable temporal
\texttt{Time} como clave de unión, una vez normalizada la columna de tiempo en todos
los ficheros. La tabla~\ref{tab:datasets} resume los datasets y su función.

\begin{table}[H]
\centering
\caption{Datasets integrados en el Sprint 2}
\label{tab:datasets}
\begin{tabular}{lll}
\toprule
\textbf{Identificador} & \textbf{Tipo de variable} & \textbf{Variables principales} \\
\midrule
\texttt{air\_temp}   & Microclima atmosférico & Temperatura del aire (WS, S1, S2) \\
\texttt{epar}        & Cultivo                & Radiación fotosintética (S1, S2) \\
\texttt{irradiance}  & Energético             & GPOA, Albedo (S1, S2) \\
\texttt{pv\_temp}    & Energético             & Temperatura de panel (S1, S2) \\
\texttt{soil\_temp}  & Cultivo                & Temperatura del suelo (S1, S2) \\
\texttt{soil\_vwc}   & Cultivo                & Contenido volumétrico de agua \\
\texttt{tracking}    & Control operativo      & Ángulo de rotación (M01--M09) \\
\texttt{wind\_speed} & Meteorológico          & Velocidad del viento \\
\bottomrule
\end{tabular}
\end{table}

\subsection{Funciones del pipeline}

Se formalizan tres funciones reutilizables dentro del notebook, que extienden el pipeline
del Sprint~1:

\begin{description}[style=nextline]
    \item[\texttt{strip\_unit(series)}]
        Extrae el valor numérico de columnas que contienen unidades embebidas en el
        texto (por ejemplo, \texttt{"41.4°"} o \texttt{"1024 W/m²"}). En Sprint~1
        solo funcionaba con columnas de tipo \texttt{object}; en Sprint~2 se extiende
        para soportar también el tipo \texttt{string} introducido en \texttt{pandas} 3.x.

    \item[\texttt{load\_csv(filepath)}]
        Carga un CSV corrigiendo la cabecera BOM UTF-8 y las líneas auxiliares de tipo
        \texttt{sep=} que generan algunos exportadores. Normaliza el nombre de la columna
        temporal a \texttt{Time} y la convierte a \texttt{datetime}.

    \item[\texttt{resample\_to\_6h(df)}]
        Reagrega cualquier dataframe numérico al grid temporal de seis horas mediante
        la media aritmética (\texttt{mean}). Esto permite unificar las distintas
        frecuencias de muestreo de los sensores.
\end{description}

\subsection{Cobertura temporal}

Una vez integrados los ocho datasets y resampleados a seis horas, el dataset combinado
contiene \textbf{1.462 filas} y \textbf{30 variables} base. La cobertura de datos no
es uniforme entre sensores; la tabla~\ref{tab:cobertura} muestra el porcentaje de
registros con valor disponible para cada grupo de variables.

\begin{table}[H]
\centering
\caption{Cobertura de datos por grupo de variables (\% de timesteps con dato válido)}
\label{tab:cobertura}
\begin{tabular}{lc}
\toprule
\textbf{Grupo de variables} & \textbf{Cobertura (\%)} \\
\midrule
Irradiancia (GPOA, Albedo, Tpanel)  & 28.5 \\
Velocidad de viento                  & 54.5 \\
VWC del suelo                        & 58.8 \\
Temperatura del suelo                & 61.5 \\
Temperatura del aire, ePAR           & 64.2 \\
Tracking (ángulo de trackers)        & 69.5 \\
\bottomrule
\end{tabular}
\end{table}

La baja cobertura de irradiancia (28.5\%) y de viento (54.5\%) limita la ventana de
solapamiento total útil para modelización. La figura~\ref{fig:cobertura} muestra el
mapa de cobertura temporal completo.

\begin{figure}[H]
    \centering
    \includegraphics[width=\textwidth]{outputs_sprint2/dataset_integrado_cobertura.png}
    \caption{Mapa de cobertura temporal del dataset integrado. Cada celda representa
             la disponibilidad de dato (blanco = disponible, negro = ausente).}
    \label{fig:cobertura}
\end{figure}

% ─── 3. DIAGNÓSTICO DE TRACKERS ─────────────────────────────────────────────
\section{Diagnóstico Previo: Estado de los Trackers}

Antes de proceder a la modelización, se realizó un diagnóstico de varianza sobre todos
los trackers disponibles. El objetivo era detectar posibles trackers en posición
de \textit{stow} fija (sin movimiento real durante el período), que en Sprint~1 se
habían marcado como sospechosos para los trackers M02, M06 y M10.

El análisis de varianza del ángulo registrado no muestra evidencia concluyente de
\textit{stow} fijo para ninguno de los tres trackers mencionados. Esto significa dos
cosas: (1) no hay suficiente evidencia para declarar fallo mecánico solo con este
análisis, y (2) la validación con el equipo de planta sigue siendo necesaria.

Se han usado únicamente los trackers con varianza confirmada (M01, M03, M05, M07, M09)
para calcular el ángulo medio representativo del sistema (\texttt{track\_mean}).

\begin{figure}[H]
    \centering
    \includegraphics[width=0.85\textwidth]{outputs_sprint2/tracker_variance_diagnostic.png}
    \caption{Varianza del ángulo registrado por tracker. La línea roja discontinua
             indica el umbral de baja varianza (5 grados$^2$) empleado como señal
             de alerta de posible \textit{stow} fijo.}
    \label{fig:tracker_var}
\end{figure}

% ─── 4. FEATURE ENGINEERING ─────────────────────────────────────────────────
\newpage
\section{Feature Engineering}

\subsection{Variables construidas}

A partir del dataset integrado se construyen las siguientes variables derivadas para
alimentar el modelo:

\subsubsection{Variables temporales y astronómicas}

\begin{description}[noitemsep, style=nextline]
    \item[\texttt{hour\_of\_day}] Hora del timestep (valores: 0, 6, 12, 18 al trabajar
        a resolución de 6h).
    \item[\texttt{day\_of\_year}] Día del año (1--365), como proxy de la estación.
    \item[\texttt{solar\_elevation\_deg}] Elevación solar teórica en grados, calculada
        mediante la fórmula astronómica estándar usando latitud $41.39°$N como proxy
        de la localización de la planta (pendiente de confirmar con el equipo técnico).
        Los valores negativos (nocturnos) se recortan a cero.
\end{description}

\subsubsection{Variables de régimen de tracking}

\begin{description}[noitemsep, style=nextline]
    \item[\texttt{track\_mean}] Media del ángulo registrado por los trackers activos
        M01, M03, M05, M07 y M09.
    \item[\texttt{tracking\_regime}] Clasificación operativa del modo de operación:
        \begin{itemize}[noitemsep]
            \item \textbf{TRACKING\_PM}: ángulo medio positivo ($> 5°$), paneles
                girados hacia el oeste (tarde).
            \item \textbf{TRACKING\_AM}: ángulo medio negativo ($< -5°$), paneles
                girados hacia el este (mañana).
            \item \textbf{HORIZONTAL}: ángulo cercano a cero ($|ángulo| < 5°$).
            \item \textbf{STOW}: ángulo próximo a 50.6° (posición de seguridad).
            \item \textbf{UNKNOWN}: ángulo no disponible.
        \end{itemize}
\end{description}

\subsubsection{Variables agregadas por sección}

Para reducir la dimensionalidad y capturar el comportamiento conjunto de cada zona de
la planta, se calculan medias por sección:

\begin{itemize}[noitemsep]
    \item \texttt{ePAR\_S1\_mean}, \texttt{ePAR\_S2\_mean}: media de los sensores ePAR
        de cada sección.
    \item \texttt{VWC\_S1\_mean}, \texttt{VWC\_S2\_mean}: media del contenido volumétrico
        de agua del suelo.
    \item \texttt{Tsoil\_S1\_mean}, \texttt{Tsoil\_S2\_mean}: media de temperatura del
        suelo.
\end{itemize}

\subsubsection{Variables con efecto retardado (lag)}

Para capturar la inercia térmica del suelo, se incorporan dos variables de lag:

\begin{itemize}[noitemsep]
    \item \texttt{Tsoil\_S1\_mean\_lag\_6h}: temperatura media del suelo de la sección~1
        seis horas antes.
    \item \texttt{Tsoil\_S1\_mean\_lag\_12h}: temperatura media del suelo de la sección~1
        doce horas antes.
\end{itemize}

\subsubsection{Variables diferenciales}

\begin{itemize}[noitemsep]
    \item \texttt{VWC\_diff\_S1\_minus\_S2}: diferencia de contenido de agua entre sección~1
        y sección~2, como indicador de heterogeneidad hídrica entre zonas.
\end{itemize}

\subsection{Nota sobre la exclusión de GPOA}

Las variables \texttt{GPOA\_S1} y \texttt{GPOA\_S2} (irradiancia sobre plano inclinado)
se calculan y se usan para construir el target IEC (componente energética), pero se
excluyen \textbf{deliberadamente} de las features del modelo. Incluirlas crearía una
circularidad directa: el modelo aprendería a predecir el IEC usando las mismas variables
irradiativas que lo definen, produciendo un R$^2$ artificialmente elevado sin valor
informativo real. Como proxies de las condiciones radiativas se conservan
\texttt{Albedo\_S1} y \texttt{Albedo\_S2}, que reflejan las condiciones de luz sin
participar directamente en el cálculo del target.

% ─── 5. DEFINICIÓN DEL TARGET: IEC ──────────────────────────────────────────
\newpage
\section{Definición del Target: Índice Energía--Cultivo (IEC)}

\subsection{Justificación}

El proyecto SAMO no tiene como objetivo maximizar la energía producida, sino encontrar
un equilibrio dinámico entre eficiencia energética y condiciones agronómicas favorables.
Para operacionalizar este objetivo en un modelo supervisado, es necesario definir una
variable objetivo (\textit{target}) única que capture ambas dimensiones.

Dado que en esta fase no existe aún un target validado por el equipo técnico y agrícola,
se propone el \textbf{Índice Energía--Cultivo (IEC)} como formulación provisional y
explícitamente revisable.

\subsection{Formulación}

El IEC se construye como la media ponderada de dos sub-índices normalizados:

\begin{equation}
    \text{IEC} = 0.5 \cdot \text{energy\_score} + 0.5 \cdot \text{crop\_score}
    \label{eq:iec}
\end{equation}

\subsubsection{Componente energética}

\begin{equation}
    \text{energy\_score} = \min\!\left(1,\; \frac{\overline{\text{GPOA}}}{P_{95}(\overline{\text{GPOA}})}\right)
\end{equation}

donde $\overline{\text{GPOA}}$ es la media de \texttt{GPOA\_S1} y \texttt{GPOA\_S2},
y $P_{95}$ es el percentil 95 del período completo. Un \texttt{energy\_score} de 1
indica irradiancia máxima (óptimo energético); de 0 indica ausencia de irradiancia
(noche o nublado).

\subsubsection{Componente agrícola}

\begin{equation}
    \text{crop\_score} = 0.4 \cdot \text{par\_score} + 0.4 \cdot \text{vwc\_score} + 0.2 \cdot \text{tsoil\_score}
\end{equation}

donde cada sub-score se define como la proximidad normalizada a un valor objetivo:

\begin{align}
    \text{par\_score}   &= \max\!\left(0,\; 1 - \frac{|\overline{\text{ePAR}} - 600|}{600}\right) \\[4pt]
    \text{vwc\_score}   &= \max\!\left(0,\; 1 - \frac{|\overline{\text{VWC}} - 30|}{15}\right) \\[4pt]
    \text{tsoil\_score} &= \max\!\left(0,\; 1 - \frac{|\overline{T_{suelo}} - 28|}{10}\right)
\end{align}

Los valores de referencia (ePAR = 600 $\mu$mol/m$^2$/s, VWC = 30\%, $T_{suelo}$ = 28°C)
son supuestos operativos provisionales que deben ser validados con el especialista
agrícola y ajustados al tipo de cultivo específico de la planta.

\subsection{Propiedades del IEC}

El IEC toma valores en el intervalo $[0, 1]$ donde:
\begin{itemize}[noitemsep]
    \item \textbf{IEC próximo a 1}: situación de alto equilibrio entre producción
        energética y condiciones agronómicas favorables.
    \item \textbf{IEC próximo a 0}: situación desfavorable en ambas dimensiones
        (noche, sequía, estrés térmico extremo, etc.).
    \item Los pesos (0.5/0.5 entre energía y cultivo; 0.4/0.4/0.2 entre componentes
        agrícolas) son configurables y no han sido validados con criterio agronómico.
\end{itemize}

\subsection{Distribución del IEC}

La tabla~\ref{tab:iec_summary} muestra los estadísticos descriptivos del IEC sobre
las 337 observaciones modelables.

\begin{table}[H]
\centering
\caption{Estadísticos descriptivos del IEC y sus componentes (n = 337)}
\label{tab:iec_summary}
\begin{tabular}{lcccc}
\toprule
\textbf{Estadístico} & \textbf{energy\_score} & \textbf{crop\_score} & \textbf{IEC} \\
\midrule
count  & 416    & 860    & 337   \\
media  & 0.339  & 0.392  & 0.397 \\
desv.\ est. & 0.349 & 0.165 & 0.248 \\
mín.   & 0.000  & 0.112  & 0.063 \\
Q25    & 0.000  & 0.251  & 0.183 \\
mediana & 0.225 & 0.378  & 0.360 \\
Q75    & 0.599  & 0.517  & 0.598 \\
máx.   & 1.000  & 0.940  & 0.925 \\
\bottomrule
\end{tabular}
\end{table}

El IEC muestra una distribución asimétrica positiva, con mediana de 0.36 y máximo de
0.925. La alta varianza del \texttt{energy\_score} (std = 0.35) refleja la alternancia
entre períodos diurnos con alta irradiancia y períodos nocturnos o nublados. La
figura~\ref{fig:iec} muestra la evolución temporal y la distribución del IEC.

\begin{figure}[H]
    \centering
    \includegraphics[width=\textwidth]{outputs_sprint2/iec_overview.png}
    \caption{Izquierda: evolución temporal del \texttt{energy\_score},
             \texttt{crop\_score} e IEC sobre el período modelable (jun--oct 2025).
             Derecha: distribución del IEC (histograma, 35 bins).}
    \label{fig:iec}
\end{figure}

% ─── 6. DATASET DE MODELIZACIÓN ─────────────────────────────────────────────
\newpage
\section{Dataset de Modelización}

\subsection{Selección del período útil}

El dataset de modelización se construye a partir de las observaciones que cumplen
simultáneamente:
\begin{enumerate}[noitemsep]
    \item El IEC está calculado (no es nulo).
    \item Al menos el 75\% de las features tienen valor disponible.
\end{enumerate}

Aplicando estos criterios sobre el dataset integrado (1.462 filas), se obtienen
\textbf{337 observaciones modelables}, correspondientes al período
\textbf{2025-06-18 -- 2025-10-05}. Este recorte es consecuencia directa de la baja
cobertura de irradiancia (28.5\%), que es condición necesaria para calcular el
\texttt{energy\_score}.

\subsection{Partición temporal}

Se utiliza una partición temporal 80/20, respetando el orden cronológico de las
observaciones para evitar fuga de información del futuro al pasado (\textit{data leakage}
temporal):

\begin{table}[H]
\centering
\caption{Partición del dataset de modelización}
\label{tab:split}
\begin{tabular}{lccc}
\toprule
\textbf{Subconjunto} & \textbf{Filas} & \textbf{Inicio} & \textbf{Fin} \\
\midrule
Train & 269 & 2025-06-18 12:00 & 2025-09-07 18:00 \\
Test  &  68 & 2025-09-08 00:00 & 2025-10-05 00:00 \\
\midrule
\textbf{Total} & \textbf{337} & \textbf{2025-06-18 12:00} & \textbf{2025-10-05 00:00} \\
\bottomrule
\end{tabular}
\end{table}

\subsection{Features del modelo}

El vector de features final contiene \textbf{19 variables}, excluyendo explícitamente
\texttt{GPOA\_S1} y \texttt{GPOA\_S2} por circularidad con el target (ver
sección~4.2). La tabla~\ref{tab:features} lista las features empleadas, agrupadas
por categoría.

\begin{table}[H]
\centering
\caption{Features del modelo y categoría}
\label{tab:features}
\begin{tabular}{ll}
\toprule
\textbf{Categoría} & \textbf{Features} \\
\midrule
Temporal / astronómica
    & \texttt{hour\_of\_day}, \texttt{day\_of\_year}, \texttt{solar\_elevation\_deg} \\[4pt]
Control operativo
    & \texttt{track\_mean}, \texttt{tracking\_regime} \\[4pt]
Condiciones radiativas (proxy)
    & \texttt{Albedo\_S1}, \texttt{Albedo\_S2} \\[4pt]
Temperatura del aire
    & \texttt{Tair\_WS}, \texttt{Tair\_S1\_center} \\[4pt]
Viento
    & \texttt{wind\_speed\_kmh} \\[4pt]
Suelo -- temperatura
    & \texttt{Tsoil\_S1\_mean}, \texttt{Tsoil\_S2\_mean},
      \texttt{Tsoil\_S1\_mean\_lag\_6h}, \texttt{Tsoil\_S1\_mean\_lag\_12h} \\[4pt]
Suelo -- humedad
    & \texttt{VWC\_S1\_mean}, \texttt{VWC\_S2\_mean}, \texttt{VWC\_diff\_S1\_minus\_S2} \\[4pt]
Cultivo
    & \texttt{ePAR\_S1\_mean}, \texttt{ePAR\_S2\_mean} \\
\bottomrule
\end{tabular}
\end{table}

% ─── 7. MODELOS ENTRENADOS ───────────────────────────────────────────────────
\newpage
\section{Modelos Entrenados}

Se han entrenado dos modelos con objetivos complementarios: uno para
interpretabilidad (extracción de reglas operativas) y otro para generalización
predictiva.

\subsection{Decision Tree Regressor}

El árbol de decisión se entrena con el objetivo de obtener reglas operativas del tipo
\textit{si--entonces} que puedan trasladarse directamente a lógica de control. La
configuración prioriza la robustez frente al sobreajuste:

\begin{table}[H]
\centering
\caption{Hiperparámetros del Decision Tree Regressor}
\begin{tabular}{ll}
\toprule
\textbf{Parámetro} & \textbf{Valor} \\
\midrule
\texttt{max\_depth}        & 4 \\
\texttt{min\_samples\_leaf} & 20 \\
\texttt{random\_state}     & 42 \\
\bottomrule
\end{tabular}
\end{table}

El preprocesado aplica imputación por mediana para variables numéricas y
\textit{one-hot encoding} para \texttt{tracking\_regime}.

\subsection{ElasticNetCV}

La regresión regularizada ElasticNet combina penalización L1 (selección de variables)
y L2 (estabilidad). Se usa validación cruzada de 5 particiones para seleccionar
automáticamente la ratio L1/L2 (\texttt{l1\_ratio}) y la intensidad de regularización
($\alpha$).

\begin{table}[H]
\centering
\caption{Hiperparámetros del ElasticNetCV}
\begin{tabular}{ll}
\toprule
\textbf{Parámetro} & \textbf{Valor} \\
\midrule
\texttt{l1\_ratio}   & [0.1, 0.5, 0.7, 0.9, 1.0] (búsqueda CV) \\
\texttt{alphas}      & 30 valores en $[10^{-3}, 10^{1}]$ (escala logarítmica) \\
\texttt{cv}          & 5 particiones temporales \\
\texttt{max\_iter}   & 20.000 \\
\bottomrule
\end{tabular}
\end{table}

El preprocesado aplica imputación por mediana y estandarización (z-score) para
variables numéricas.

% ─── 8. RESULTADOS ──────────────────────────────────────────────────────────
\newpage
\section{Resultados de Evaluación}

\subsection{Métricas sobre el conjunto de test}

La tabla~\ref{tab:metricas} muestra las métricas de evaluación sobre el conjunto de
test temporal (n = 68, período septiembre--octubre 2025).

\begin{table}[H]
\centering
\caption{Métricas de evaluación en el conjunto de test temporal}
\label{tab:metricas}
\begin{tabular}{lccc}
\toprule
\textbf{Modelo} & \textbf{MAE} & \textbf{RMSE} & \textbf{R$^2$} \\
\midrule
ElasticNetCV       & 0.0572 & 0.0803 & 0.873 \\
Decision Tree      & 0.0912 & 0.1101 & 0.762 \\
\bottomrule
\end{tabular}
\end{table}

\subsubsection{Lectura de resultados}

\begin{itemize}
    \item \textbf{ElasticNet} generaliza mejor, con un R$^2$ de 0.873. El error
        absoluto medio de 0.057 puntos sobre una escala $[0, 1]$ indica una capacidad
        predictiva sólida para el período de test.

    \item \textbf{Decision Tree} rinde peor en términos predictivos (R$^2$ = 0.762),
        pero su ventaja reside en la interpretabilidad directa de las reglas generadas.
        Un R$^2$ de 0.76 sigue siendo suficientemente bueno para extraer reglas
        operativas con soporte estadístico.

    \item Los valores de R$^2$ son más modestos que en versiones previas del modelo
        donde \texttt{GPOA} se incluía como feature (0.95 para ElasticNet). Esa
        diferencia es \textbf{esperada y correcta}: la versión anterior producía un
        resultado artificialmente inflado por circularidad con el target.
\end{itemize}

\begin{figure}[H]
    \centering
    \includegraphics[width=\textwidth]{outputs_sprint2/model_predictions_vs_actual.png}
    \caption{Comparación entre valores reales y predichos de IEC sobre el conjunto de
             test para ambos modelos. La línea discontinua negra representa la
             predicción perfecta.}
    \label{fig:predictions}
\end{figure}

% ─── 9. INTERPRETABILIDAD ───────────────────────────────────────────────────
\newpage
\section{Interpretabilidad de los Modelos}

\subsection{Importancia de variables en el árbol}

La tabla~\ref{tab:tree_importance} muestra la importancia relativa de cada variable
en el árbol de decisión. El árbol concentra casi toda su capacidad explicativa en
cuatro variables.

\begin{table}[H]
\centering
\caption{Importancia de variables en el Decision Tree (top 6)}
\label{tab:tree_importance}
\begin{tabular}{lcc}
\toprule
\textbf{Variable} & \textbf{Importancia} & \textbf{\%} \\
\midrule
\texttt{Albedo\_S1}          & 0.773 & 77.3\% \\
\texttt{Albedo\_S2}          & 0.112 & 11.2\% \\
\texttt{solar\_elevation\_deg} & 0.066 & 6.6\% \\
\texttt{VWC\_S1\_mean}       & 0.048 & 4.8\% \\
\texttt{Tair\_S1\_center}    & 0.001 & 0.1\% \\
resto (14 variables)         & 0.000 & 0.0\% \\
\bottomrule
\end{tabular}
\end{table}

El dominio de \texttt{Albedo\_S1} (77.3\%) refleja que las condiciones radiativas son
el principal determinante del IEC en este período. Esto es coherente con el hecho de
que la componente energética representa el 50\% del IEC y está directamente relacionada
con la irradiancia. \texttt{tracking\_regime} y \texttt{track\_mean} no aparecen en el
árbol, lo que indica que, una vez fijadas las condiciones meteorológicas, el árbol no
discrimina entre regímenes de rotación.

\begin{figure}[H]
    \centering
    \includegraphics[width=\textwidth]{outputs_sprint2/decision_tree_interpretability.png}
    \caption{Izquierda: top 10 variables por importancia en el árbol. Derecha:
             representación visual del árbol de decisión (profundidad máxima 4).}
    \label{fig:tree}
\end{figure}

\subsection{Estructura del árbol de decisión}

El árbol se puede describir textualmente con las siguientes ramas principales:

\begin{itemize}
    \item \textbf{Condición raíz}: Si \texttt{Albedo\_S1} $> 55.7$ → rama de alta irradiancia.
    \item \textbf{Rama alta irradiancia}:
    \begin{itemize}[noitemsep]
        \item Si \texttt{solar\_elevation\_deg} $> 68°$ → IEC esperado = \textbf{0.862}
            (mejor escenario).
        \item Si \texttt{solar\_elevation\_deg} $\in (54.2°, 68°]$ → IEC = \textbf{0.692}.
        \item Si \texttt{solar\_elevation\_deg} $\leq 54.2°$ → IEC = \textbf{0.583}.
    \end{itemize}
    \item \textbf{Rama baja irradiancia} (\texttt{Albedo\_S1} $\leq 55.7$):
    \begin{itemize}[noitemsep]
        \item Si \texttt{Albedo\_S2} $> 18.0$ y \texttt{VWC\_S1} $> 20.0$ → IEC = \textbf{0.435}.
        \item Si \texttt{Albedo\_S2} $\leq 18.0$ y \texttt{VWC\_S1} $> 22.2$ → IEC $\approx$ \textbf{0.276}.
        \item Si \texttt{Albedo\_S2} $\leq 18.0$ y \texttt{VWC\_S1} $\leq 22.2$ → IEC $\approx$ \textbf{0.139}
            (peor escenario).
    \end{itemize}
\end{itemize}

\subsection{Coeficientes del ElasticNet}

La tabla~\ref{tab:coef} muestra los coeficientes estandarizados más influyentes del
modelo ElasticNet. Los signos indican la dirección del efecto sobre el IEC.

\begin{table}[H]
\centering
\caption{Coeficientes estandarizados del ElasticNet (top 10 por magnitud absoluta)}
\label{tab:coef}
\begin{tabular}{lcc}
\toprule
\textbf{Variable} & \textbf{Coeficiente} & \textbf{Dirección} \\
\midrule
\texttt{Albedo\_S2}                      & +0.123 & Positivo \\
\texttt{tracking\_regime\_HORIZONTAL}    & --0.114 & Negativo \\
\texttt{tracking\_regime\_TRACKING\_PM}  & +0.071 & Positivo \\
\texttt{Tsoil\_S2\_mean}                 & +0.046 & Positivo \\
\texttt{VWC\_S2\_mean}                   & +0.036 & Positivo \\
\texttt{Tsoil\_S1\_mean}                 & --0.035 & Negativo \\
\texttt{Tair\_WS}                        & +0.034 & Positivo \\
\texttt{Tair\_S1\_center}                & --0.031 & Negativo \\
\texttt{VWC\_S1\_mean}                   & +0.025 & Positivo \\
\texttt{day\_of\_year}                   & --0.024 & Negativo \\
\bottomrule
\end{tabular}
\end{table}

\textbf{Interpretación de los coeficientes más relevantes:}

\begin{itemize}
    \item \texttt{tracking\_regime\_HORIZONTAL} tiene el segundo coeficiente más alto en
        valor absoluto y es \textbf{negativo} (--0.114). Esto confirma que operar en
        modo horizontal penaliza el IEC respecto a otros modos de seguimiento.

    \item \texttt{tracking\_regime\_TRACKING\_PM} tiene coeficiente \textbf{positivo}
        (+0.071), lo que indica que el seguimiento solar de tarde está asociado a mejor
        IEC.

    \item \texttt{Albedo\_S2} positivo (+0.123) refleja que mayores condiciones
        radiativas en la sección~2 contribuyen positivamente al equilibrio
        energía--cultivo.

    \item Las variables \texttt{Tsoil} muestran efectos opuestos entre secciones, lo que
        sugiere una dinámica térmica diferencial entre S1 y S2 que merece análisis
        adicional.
\end{itemize}

\begin{figure}[H]
    \centering
    \includegraphics[width=0.85\textwidth]{outputs_sprint2/elasticnet_coefficients.png}
    \caption{Coeficientes estandarizados del ElasticNet. Verde: efecto positivo sobre
             el IEC. Naranja: efecto negativo. Solo se muestran los 12 más influyentes.}
    \label{fig:elastic}
\end{figure}

% ─── 10. ANÁLISIS POR RÉGIMEN ────────────────────────────────────────────────
\newpage
\section{Análisis por Régimen de Operación}

\subsection{Comparativa de IEC por régimen}

El análisis observacional directo por régimen de tracking es la base principal para
las reglas candidatas de rotación. La tabla~\ref{tab:regime} muestra el IEC estadístico
por régimen para las 337 observaciones modelables.

\begin{table}[H]
\centering
\caption{IEC descriptivo por régimen de operación}
\label{tab:regime}
\begin{tabular}{lccc}
\toprule
\textbf{Régimen} & \textbf{N observaciones} & \textbf{IEC medio} & \textbf{IEC mediana} \\
\midrule
TRACKING\_PM & 84  & \textbf{0.742} & \textbf{0.741} \\
TRACKING\_AM & 84  & 0.440          & 0.429 \\
HORIZONTAL   & 169 & 0.204          & 0.184 \\
\bottomrule
\end{tabular}
\end{table}

Los resultados muestran una diferencia muy marcada entre regímenes:

\begin{itemize}
    \item \textbf{TRACKING\_PM} supera a HORIZONTAL en un \textbf{+0.538 puntos de IEC}
        de media (equivalente a un 264\% de mejora relativa).
    \item \textbf{TRACKING\_AM} supera a HORIZONTAL en +0.236 puntos.
    \item \textbf{HORIZONTAL} queda como el régimen de peor equilibrio energía--cultivo
        en las condiciones observadas durante junio--octubre 2025.
\end{itemize}

\begin{figure}[H]
    \centering
    \includegraphics[width=\textwidth]{outputs_sprint2/policy_rules_diagnostics.png}
    \caption{Izquierda: IEC medio por franja horaria y régimen (mapa de calor).
             Derecha: distribución del IEC por régimen (diagrama de caja).}
    \label{fig:policy}
\end{figure}

\subsection{Casos de IEC alto por franja horaria}

El análisis de casos con IEC superior al percentil 75 ($\text{IEC} \geq 0.598$) por
franja horaria y régimen revela los dos escenarios de operación más favorables:

\begin{table}[H]
\centering
\caption{Franjas horarias con IEC alto (percentil 75) por régimen}
\label{tab:high_iec}
\begin{tabular}{lccccc}
\toprule
\textbf{Hora} & \textbf{Régimen} & \textbf{N} & \textbf{IEC mediana}
    & \textbf{Ángulo mediano} & \textbf{VWC\_S1 mediana} \\
\midrule
12:00 & TRACKING\_PM & 75 & 0.769 & +31.8° & 24.0\% \\
06:00 & TRACKING\_AM & 10 & 0.602 & --32.2° & 29.2\% \\
\bottomrule
\end{tabular}
\end{table}

La franja del mediodía con seguimiento PM es el escenario de mayor IEC observado
con soporte estadístico significativo (n = 75). La franja de las 06:00 con seguimiento
AM tiene un soporte mucho más reducido (n = 10) y debe interpretarse con cautela.

% ─── 11. REGLAS CANDIDATAS DE ROTACIÓN ──────────────────────────────────────
\newpage
\section{Reglas Candidatas de Política de Rotación}

\subsection{Método de derivación}

Las reglas candidatas se derivan de la combinación de tres fuentes complementarias:
(1) los nodos del árbol de decisión, (2) el análisis observacional por régimen
(tabla~\ref{tab:regime}), y (3) los casos de IEC alto por franja horaria
(tabla~\ref{tab:high_iec}). Las reglas son de carácter \textit{observacional} y
correlacional, no causal: describen qué configuraciones coincidieron con mejor IEC
en el período junio--octubre 2025, pero no implican causalidad probada.

\subsection{Las seis reglas candidatas}

A continuación se presentan las seis reglas derivadas, ordenadas por IEC esperado
de mayor a menor.

\vspace{0.4cm}
\noindent\textbf{Regla 1 -- Prioridad alta: condiciones de máxima productividad}

\begin{reglabox}
\textbf{Condición:} \texttt{Albedo\_S1} $>$ 55.7 Y elevación solar $>$ 68° \\
\textbf{Acción:} Priorizar tracking activo de tarde; evitar modos conservadores \\
\textbf{IEC esperado:} 0.862 \quad \textbf{Soporte:} 33 observaciones \\
\textit{Nodo de mayor IEC del árbol de decisión.}
\end{reglabox}

\vspace{0.3cm}
\noindent\textbf{Regla 2 -- Operativa: mediodía con seguimiento PM}

\begin{reglabox}
\textbf{Condición:} Franja horaria 12:00 Y sistema en TRACKING\_PM \\
\textbf{Acción:} Mantener ángulo cercano a +31.8° \\
\textbf{IEC esperado:} 0.769 \quad \textbf{Soporte:} 75 observaciones \\
\textit{Mayor soporte estadístico de todas las reglas operativas.}
\end{reglabox}

\vspace{0.3cm}
\noindent\textbf{Regla 3 -- Operativa: mañana con seguimiento AM}

\begin{reglabox}
\textbf{Condición:} Franja horaria 06:00 Y sistema en TRACKING\_AM \\
\textbf{Acción:} Mantener ángulo cercano a --32.2° \\
\textbf{IEC esperado:} 0.602 \quad \textbf{Soporte:} 10 observaciones \\
\textit{Soporte reducido; requiere validación con más datos.}
\end{reglabox}

\vspace{0.3cm}
\noindent\textbf{Regla 4 -- Transición: condiciones intermedias}

\begin{reglabox}
\textbf{Condición:} \texttt{Albedo\_S1} $\leq$ 55.7 Y \texttt{Albedo\_S2} $>$ 18.0
    Y \texttt{VWC\_S1} $>$ 20.0 \\
\textbf{Acción:} Usar tracking suave o transición controlada; evitar HORIZONTAL permanente \\
\textbf{IEC esperado:} 0.435 \quad \textbf{Soporte:} 51 observaciones \\
\textit{Escenario intermedio del árbol; mejor que HORIZONTAL (IEC medio 0.204).}
\end{reglabox}

\vspace{0.3cm}
\noindent\textbf{Regla 5 -- Evitación: HORIZONTAL en franjas productivas}

\begin{reglabox}
\textbf{Condición:} Sistema en modo HORIZONTAL durante franjas diurnas productivas \\
\textbf{Acción:} Evitar este modo cuando existan condiciones para tracking activo \\
\textbf{IEC esperado:} 0.184 \quad \textbf{Soporte:} 169 observaciones \\
\textit{IEC medio HORIZONTAL = 0.204 vs. TRACKING\_PM = 0.742.
    El diferencial es de +0.538 puntos.}
\end{reglabox}

\vspace{0.3cm}
\noindent\textbf{Regla 6 -- Cautela hídrica: condiciones de excepción}

\begin{reglabox}
\textbf{Condición:} \texttt{Albedo\_S1} $\leq$ 55.7 Y \texttt{Albedo\_S2} $\leq$ 18.0
    Y \texttt{VWC\_S1} $\leq$ 22.2 Y \texttt{GPOA\_S2} $\leq$ 35.5 \\
\textbf{Acción:} No forzar política orientada a producción; tratar como escenario
    de excepción \\
\textbf{IEC esperado:} 0.093 \quad \textbf{Soporte:} 43 observaciones \\
\textit{Nodo de peor IEC del árbol. Baja irradiancia y baja humedad de suelo
    simultáneas.}
\end{reglabox}

\subsection{Resumen operativo}

Si hubiera que traducir las reglas a una sola idea operativa:

\begin{center}
\fbox{\parbox{0.88\textwidth}{
    \centering\vspace{0.2cm}
    En las horas productivas, operar con tracking activo y ángulos intermedios
    (~$-$32° por la mañana y ~+32° al mediodía) produce un IEC hasta 3.6 veces
    mayor que operar en posición horizontal. Las reglas de cautela se reservan
    para escenarios de baja irradiancia simultánea con baja humedad del suelo.
    \vspace{0.2cm}
}}
\end{center}

% ─── 12. LIMITACIONES ───────────────────────────────────────────────────────
\newpage
\section{Limitaciones y Trabajo Futuro}

\subsection{Limitaciones del Sprint 2}

\begin{enumerate}
    \item \textbf{IEC provisional}: Los pesos del IEC (0.5/0.5 entre energía y cultivo;
        0.4/0.4/0.2 entre componentes agrícolas) y los valores de referencia (ePAR = 600,
        VWC = 30\%, $T_{suelo}$ = 28°C) son supuestos operativos sin validar con
        el especialista agrícola ni con el tipo de cultivo específico de la planta.

    \item \textbf{Ventana temporal parcial}: El período modelable (junio--octubre 2025)
        corresponde a los meses de verano y principio de otoño. El modelo no ha visto
        datos de invierno ni primavera, lo que limita la generalización de las reglas
        a otras estaciones.

    \item \textbf{N reducido}: 337 observaciones modelables con 68 en test es un
        tamaño de muestra pequeño para evaluación robusta. Los resultados deben
        interpretarse como indicativos, no como estimaciones definitivas.

    \item \textbf{Árbol sin uso de variables de control}: El árbol de decisión no usa
        \texttt{tracking\_regime} ni \texttt{track\_mean} como variables de split. Esto
        refleja que, en este período, las condiciones meteorológicas dominan la
        variabilidad del IEC y el árbol las captura antes de llegar a las variables de
        control. Las reglas operativas de rotación se derivan del análisis observacional
        por régimen, no del árbol directamente.

    \item \textbf{Correlación vs.\ causalidad}: Las reglas son correlacionales. El
        sistema observó que ciertos regímenes coincidieron con mejor IEC, pero no se
        puede concluir causalidad sin un diseño experimental controlado. La validación
        con el especialista agrícola es imprescindible antes de automatizar estas reglas.

    \item \textbf{Latitud como proxy}: El ángulo solar teórico se calcula con latitud
        41.39°N. Si la planta está en otra latitud, este valor introduce un error
        sistemático en la variable \texttt{solar\_elevation\_deg}.

    \item \textbf{Trackers M02, M06, M10}: El estado real de estos trackers no ha podido
        confirmarse con el equipo de planta. Si alguno estuviera en stow fijo permanente,
        el \texttt{track\_mean} podría estar sesgado.
\end{enumerate}

\subsection{Trabajo futuro (Sprint 3 y Sprint 4)}

\begin{itemize}[noitemsep]
    \item Validar el IEC y sus umbrales con el especialista agrícola antes de integrarlo
        en el dashboard.
    \item Integrar las 6 reglas candidatas en la lógica de recomendación del dashboard,
        marcándolas como pendientes de validación agronómica.
    \item Revisar el estado de los trackers M02, M06 y M10 con el equipo de planta.
    \item Confirmar la latitud real de la instalación para corregir \texttt{solar\_elevation\_deg}.
    \item Cuando se disponga de datos de invierno/primavera, reentrenar el modelo para
        cubrir el ciclo anual completo.
\end{itemize}

% ─── 13. CONCLUSIONES ───────────────────────────────────────────────────────
\section{Conclusiones}

El Sprint~2 ha producido una primera política de rotación cuantificada y trazable,
con los siguientes hallazgos principales:

\begin{enumerate}
    \item \textbf{El régimen TRACKING\_PM es el de mayor IEC observado}, con una
        mediana de 0.741, frente a 0.184 de HORIZONTAL. La diferencia absoluta de
        +0.557 en la mediana representa el potencial de mejora más directo y accionable
        del sistema.

    \item \textbf{Se han definido 6 reglas operativas candidatas}, con soporte
        estadístico documentado y trazabilidad a los datos de origen. Las dos reglas
        con mayor soporte (n = 75 y n = 51 respectivamente) apuntan al tracking activo
        en franjas productivas y a evitar el modo horizontal.

    \item \textbf{El modelo ElasticNet alcanza R$^2$ = 0.873} con un error absoluto medio
        de 0.057 sobre el conjunto de test temporal, confirmando la capacidad predictiva
        del IEC como variable objetivo.

    \item \textbf{El pipeline de datos queda formalizado y reproducible}: las funciones
        de limpieza, integración y resampleado están documentadas en el notebook y son
        reutilizables para el Sprint~3.

    \item \textbf{El IEC propuesto es operativo pero provisional}: su validación con
        criterio agronómico es la principal deuda técnica que debe resolverse antes
        de la entrega final.

\end{enumerate}

El Sprint~3 debe centrarse en integrar estas reglas y el modelo predictivo en el
dashboard operativo, con las alertas y visualizaciones necesarias para facilitar la
toma de decisiones al equipo de planta.

% ─── APÉNDICE ───────────────────────────────────────────────────────────────
\newpage
\appendix
\section{Archivos generados}

Los siguientes archivos se han generado en la carpeta \texttt{sprint2/outputs\_sprint2/}:

\begin{table}[H]
\centering
\caption{Outputs del Sprint 2}
\begin{tabular}{ll}
\toprule
\textbf{Archivo} & \textbf{Contenido} \\
\midrule
\texttt{dataset\_integrado\_6h.csv}          & Dataset completo integrado a 6h \\
\texttt{dataset\_integrado\_cobertura.csv}   & Cobertura por variable (\%) \\
\texttt{dataset\_integrado\_cobertura.png}   & Mapa de cobertura temporal \\
\texttt{dataset\_modelizacion\_6h.csv}       & Dataset final de modelización (337 filas) \\
\texttt{iec\_summary.csv}                    & Estadísticos descriptivos del IEC \\
\texttt{iec\_overview.png}                   & Evolución temporal e histograma del IEC \\
\texttt{tracker\_variance\_diagnostic.csv}   & Varianza por tracker \\
\texttt{tracker\_variance\_diagnostic.png}   & Gráfico diagnóstico de trackers \\
\texttt{model\_metrics.csv}                  & Métricas MAE/RMSE/R² por modelo \\
\texttt{model\_predictions\_vs\_actual.png}  & Real vs.\ predicho (test set) \\
\texttt{decision\_tree\_rules.txt}           & Árbol en texto \\
\texttt{decision\_tree\_feature\_importance.csv} & Importancia de variables \\
\texttt{decision\_tree\_interpretability.png}    & Árbol visualizado \\
\texttt{elasticnet\_coefficients.csv}        & Coeficientes del ElasticNet \\
\texttt{elasticnet\_coefficients.png}        & Gráfico de coeficientes \\
\texttt{regime\_summary\_iec.csv}            & IEC estadístico por régimen \\
\texttt{high\_iec\_policy\_table.csv}        & Casos de IEC alto por franja y régimen \\
\texttt{candidate\_rotation\_rules.csv}      & Las 6 reglas candidatas de rotación \\
\texttt{policy\_rules\_diagnostics.png}      & Mapa de calor y boxplot por régimen \\
\bottomrule
\end{tabular}
\end{table}

\section{Historial de versiones}

\begin{table}[H]
\centering
\begin{tabular}{llll}
\toprule
\textbf{Versión} & \textbf{Fecha} & \textbf{Autor} & \textbf{Cambios} \\
\midrule
0.1 & Abril 2026 & Equipo Chacorocalfarobar & Versión inicial \\
1.0 & Abril 2026 & Equipo Chacorocalfarobar & Revisión GPOA; carpeta outputs\_sprint2 \\
\bottomrule
\end{tabular}
\end{table}

\end{document}
